% !PRINTED 729
\weberpage{729}
%<<CONT
imprimitive numbers) and is characterised by the equation $\mathfrak{E}(r)
= \mathfrak{E}(-r)$.

By a \gesperrt{normal unit} of the field $H_m$ we shall now understand a
unit which is not $= \pm 1$, but for which the equation
\weq{3}{\mathfrak{E}(r)\mathfrak{E}(-r) = \pm 1}
is satisfied.

It is clear that a unit of the field $H_\mu$ can in no case satisfy
this condition. But not every unit in $H_m$ which does not belong
to the field $H_\mu$ will be a normal unit $^{1)}$.

The units $\tau_\beta$ [\S.\,199, (5)], however, do belong to the normal units.
For we have
\weq{4}{\tau_\beta = \mathrm{tg}\,\frac{5^\beta\pi}{m} = \frac{1-r_\beta}{1+r_\beta}r_\beta^{\nu},}
and if in this the substitution $(r,-r) = (r_\beta, r_{\beta+\nu'})$ is carried
out, then $\tau_\beta$ passes into
\weq{5}{\tau_{\beta+\nu'} = -\frac{1}{\tau_\beta}}
which corresponds to the relation (3).

A system of $\nu'$ normal units
\[
\mathfrak{E}_0(r), \mathfrak{E}_1(r), \ldots, \mathfrak{E}_{\nu'-1}(r)
\]
shall be called an \gesperrt{independent system} if the determinant
\weq{6}{\pm\Sigma\pm\log|\mathfrak{E}_0(r_0)|\log|\mathfrak{E}_1(r_1)|\ldots\log|\mathfrak{E}_{\nu'-1}(r_{\nu'-1})|
= L(\mathfrak{E}_0, \mathfrak{E}_1, \ldots, \mathfrak{E}_{\nu'-1})}
has a value different from zero.

The absolute value of this determinant, $L$, we call here too
the \gesperrt{regulator} of the system $\mathfrak{E}_0, \mathfrak{E}_1, \ldots, \mathfrak{E}_{\nu'-1}$.

The units $\tau_0, \tau_1, \ldots, \tau_{\nu'-1}$ form an independent system
of normal units, for their regulator
\[
L(\tau_0, \tau_1, \ldots, \tau_{\nu'-1})
\]
is precisely the determinant $T$ defined in the preceding section
[\S.\,199, (14), (15)]. And that this cannot be zero follows
immediately from the fact that it occurs as a factor in the expression for the
%CONT>>

\par\noindent{\footnotesize $^{1)}$ In the paper mentioned above (Acta
mathematica, Bd.\,8) the author designated the normal units as ``primitive
units''. This name is rejected here, because it could give occasion to the
error of supposing that every unit of the field $H_m$ which does not at the
same time belong to the field $H_\mu$ is to be reckoned among them.}
